\section{组织与分工}

本项目团队成员结构合理、分工明确，充分发挥了每位成员的专业特长，确保了开发工作的高效推进。

\subsection{团队角色与分工}

\begin{longtable}{|p{34mm}|p{48mm}|p{66mm}|}
\hline
\textbf{成员} & \textbf{角色} & \textbf{主要职责} \\
\hline
刘鲍非 & 产品经理 & 负责需求分析、产品原型设计与团队协调，确保项目目标与用户需求高度契合。\\
\hline
张晋恺 & 前端开发 & 负责系统前端架构设计、界面实现与用户体验优化，确保交互友好与美观性。\\
\hline
刘谦 & 后端开发 & 负责后端架构设计、核心业务逻辑实现与接口开发，保障系统高效稳定运行。\\
\hline
张磊 & 后端开发 & 参与后端模块开发与数据库设计，协助系统集成与性能优化。\\
\hline
李文耀 & 后端开发 & 负责后端功能实现，提升系统的丰富性与可扩展性。\\
\hline
施兴睿 & 数据开发 & 负责脑电信号数据处理、数据建模与分析，提升系统数据价值。\\
\hline
周左艺 & 数据开发 & 参与数据相关模块开发，协助数据可视化与数据质量保障。\\
\hline
杨吉祥 & 测试 & 负责测试用例设计、功能测试与缺陷跟踪，保障系统质量。\\
\hline
田红瑾 & 测试 & 参与系统测试与用户验收，推动产品持续优化。\\
\hline
\end{longtable}

\subsection{高效协作机制}

本项目采用了业界领先的协作与管理方式，极大提升了团队的沟通效率与项目执行力：

\begin{itemize}
	\item \textbf{代码协作：} 依托 GitHub Organization 进行代码托管与版本管理，充分利用分支管理、Pull Request、Code Review 等机制，确保代码质量与团队协作的高效有序。
	\item \textbf{文档与沟通：} 团队日常沟通与文档协作主要通过飞书团队进行，实时同步项目进展、需求变更与技术难题，保障信息透明与决策高效。
	\item \textbf{即时响应：} 通过 QQ 群组实现成员间的即时沟通与快速响应，遇到紧急问题可第一时间协同解决。
	\item \textbf{会议与任务跟踪：} 定期召开线上会议，采用飞书日历与任务看板进行任务分解与进度追踪，确保每一阶段目标的顺利达成。
\end{itemize}

通过上述多平台协作机制的有机结合，团队实现了高效沟通、敏捷开发与高质量交付，为脑电信号数据管理系统的成功落地提供了坚实保障。

\subsection{进度}

\begin{longtable}{|p{34mm}|p{48mm}|p{66mm}|}
\hline
\textbf{阶段} & \textbf{起始时间} & \textbf{完成时间} \\
\hline
立项、需求阶段 & 2026年4月27日 & 2026年5月10日 \\
\hline
设计阶段 & 2026年5月11日 & 2026年5月17日 \\
\hline
编码阶段 & 2026年5月18日 & 2026年5月31日 \\
\hline
测试阶段 & 2026年6月1日 & 2026年6月14日 \\
\hline
验收阶段 & 2026年6月15日 & 2026年6月22日 \\
\end{longtable}
